\section{From electromagnetic regions to physical machine preimages}
\label{sec:geometry-preimage}

The coordinates in Part II are not direct \gls{cad} dimensions.  Introduce a
forward electromagnetic map
\begin{equation}
 \boxed{\gls{sym:theta}=\gls{sym:machinemap}(\gls{sym:g})},
 \label{eq:geometry-forward-map}
\end{equation}
where \(\vect g\) can contain stack length, air-gap diameter, pole count,
turns, conductor area, slot and rotor dimensions, excitation turns, PM
dimensions, air gap, and material properties.  Analytical magnetic circuits,
reluctance networks, surrogate models, or \gls{fea} may instantiate
\(\mathcal M\) at different fidelities.

There is generally no single inverse \(\mathcal M^{-1}\).  Different slot,
winding, rotor, and material combinations can yield similar values of
\((R_s,R_{\mathrm{exc}},L_d,L_q,L_{\mathrm{exc}},\psi_{\mathrm{pm}})\).
The correct inverse object is the preimage
\begin{equation}
 \boxed{\gls{sym:gfeas}=
 \{\vect g:\mathcal M(\vect g)\in\Theta_{\mathrm{feas}}\}.}
 \label{eq:geometry-feasible-preimage}
\end{equation}
We first seek all physical candidates whose electromagnetic image can meet
the specification.  Only then do mass, cost, torque density, material use,
ripple, manufacturability, thermal behavior, and mechanical stress choose
among them.

\begin{figure}[tb]
\centering
\begin{tikzpicture}[node distance=1.25cm and 1.05cm,>=Latex,font=\small,
 box/.style={draw,rounded corners,align=center,text width=3.15cm,minimum height=.9cm,fill=blue!6},
 inv/.style={draw,rounded corners,align=center,text width=3.15cm,minimum height=.9cm,fill=orange!10}]
 \node[box] (g) {physical design $\vect g$\\geometry, winding, materials};
 \node[box,right=of g] (t) {electromagnetic $\boldsymbol\theta$\\$R,L,\psi,p,\ldots$};
 \node[box,right=of t] (c) {operating capability\\$M_{\max},P_{\min},\mathcal W$};
 \draw[->,thick] (g)--node[above,font=\scriptsize] {$\mathcal M$} (t);
 \draw[->,thick] (t)--node[above,font=\scriptsize] {$\mathcal C$} (c);
 \node[inv,below=1.35cm of c] (req) {required capability\\torque-speed envelope};
 \node[inv,left=of req] (tf) {$\Theta_{\mathrm{feas}}$\\admissible parameters};
 \node[inv,left=of tf] (gf) {$\mathcal G_{\mathrm{feas}}$\\non-unique preimage};
 \draw[->,thick,orange!80!black] (req)--node[above,font=\scriptsize] {level set} (tf);
 \draw[->,thick,orange!80!black] (tf)--node[above,font=\scriptsize] {preimage} (gf);
 \draw[dashed,->] (gf)--(g);
 \draw[dashed,->] (tf)--(t);
\end{tikzpicture}
\caption{Forward and inverse readings.  The physical-to-electromagnetic map is
many-to-one in general, so the inverse result is a feasible preimage rather
than a unique machine.}
\label{fig:forward-inverse-map}
\end{figure}


This separation is practically important.  The analytic circles and windows
are exact only in electromagnetic parameter space under their stated model.
Pulling them back through \(\mathcal M\) can bend, disconnect, or even empty
the corresponding region in geometry space.  Discrete slot/pole choices make
\(\mathcal G_{\mathrm{feas}}\) a union of families rather than one smooth
body.  Manufacturing and material bounds may remove parts of
\(\Theta_{\mathrm{feas}}\) that look attractive electromagnetically.

\paragraph{A hierarchy of tests.}
A candidate geometry can be processed from cheap to expensive:
\begin{enumerate}
 \item evaluate a low-order map \(\widehat{\mathcal M}(\vect g)\);
 \item reject it if the analytic loss-only capability fails;
 \item evaluate the voltage-aware support-eigenvalue value function at all
       mandatory operating points;
 \item retain only positive-margin candidates for nonlinear magnetic,
       thermal, mechanical, demagnetization, and ripple analysis;
 \item validate finalists with high-fidelity \gls{fea} and experiments.
\end{enumerate}
This is a pre-filter, not a replacement for field analysis.  Its benefit is
that obviously infeasible candidates do not require a current grid or a full
high-fidelity torque-speed map.  Published PMSM design procedures already
couple geometry, flux weakening, multiobjective search, and finite-element
verification \cite{dang2017,parasiliti2012}; the present value-function view
makes the inner feasibility calculation explicit and reusable.

The margin vector \((m_1,\ldots,m_N)\) should be passed to the outer optimizer,
not only a pass/fail flag.  It provides a smooth ranking away from active-set
changes, supports surrogate learning, and identifies which operating point
deserves higher-fidelity refinement.  Near a boundary, original equations and
independent residuals remain mandatory because small surrogate errors can
change feasibility.

